\scanpage{168}
\mspage{5}
% archive sheet number: 167
de $X$. Alors $\pi_1^{\unsure{(p')|(p)}}(X,\xi)$ est (en tant que \unsure{pro-$p'$-groupe|pro-fini-groupe}) de présentation finie.

\uline{Th. 1.4} (Th. de l'action modérée du groupe de monodromie locale). Soient $S$ un schéma strictement local, $\bar X$ un schéma sur $S$, $Y$ un sous-schéma fermé de $\bar X$, $X=\bar X-Y$.
\inserted{$U$ un ouvert de $S$ (resp. $\eta$ un pt de $U$), $\xi$ un pt géom. de $X|U$ (resp. de $X_\eta$). \deleted{\unsure{On suppose}} Dans le cas non resp., $\eta$ désigne le pt de $S$ sous $\xi$.}

a) $\bar X$, \inserted{$Y$} \deleted{sont} projectifs, de présentation finie, \deleted{\unsure{plats}} sur $S$, \inserted{l'ens des pts de $\bar X_s$ \unsure{où} $\bar X$ ou $Y$ n'est pas plat sur $S$ est de codim $\geq2$ dans $\bar X_s$.}

\deleted{b) L'ensemble des pts de $\unsure{X_s|Y_s}$ \unsure{de} \illegible{4} de codim $\geq2$ dans $\bar X_s$, et} $Y_s$ est rare dans $\bar X_s$.

c) $\bar X_s$ est, dans le complémentaire d'une partie de codim $\geq2$, un ``\deleted{diviseur}\inserted{schéma} à croisements normaux''.
\deleted{\unsure{Les fibres}} \deleted{\unsure{les pts de}} \deleted{$\bar X_\eta$ \unsure{est géométriquement} \illegible{2} \unsure{non lisse} \illegible{2} \unsure{la fibre est} $\geq2$ \unsure{(resp. $\bar X_\eta$)}}

% slanted
\margin{\deleted{\unsure{tout ceci} \illegible{3} \unsure{par une} \illegible{3}}}
% slanted
\margin{\unsure{Remplacer} l'ouvert $U$ ou le point $\eta$ par un schéma $T$ géom connexe sur $S$, avec hypothèse assurant que $I=\pi_1(T,\xi)$ opère sur $\pi_1^{(p')}(X_{\bar\eta},\xi)$.}

d) $\bar X|U$ est lisse et $Y|U$ diviseur à croisements normaux relatifs (resp. $\bar X_\eta$ a un ens de non-lissité de codim $\geq2$).
\deleted{De plus, $\bar X_\eta$ est \unsure{géométriquement} \illegible{2}, et $Y_\eta$ est \unsure{fini}.}
\inserted{$\bar X_{\bar\eta}$ est \deleted{géom intègre} irréductible.}

Soit $I=\pi_1(U,\xi)$ [resp. $I=\pi_1(\eta,\xi)\{=\operatorname{Gal}(\bar\eta/\eta)\}$], considérons l'hom naturel
\[
I\mathrel{\text{---}}\operatorname{Aut\,ext}\,\pi_1^{(p')}(X_{\bar\eta\deleted{\unsure{\bar K}}},\xi)
\]
\deleted{resp. $I\mathrel{\text{---}}\operatorname{Aut\,ext}\,\pi_1^{(p')}(X_{\bar\eta},\xi)$}

\diagram{\begin{tikzcd}[ampersand replacement=\&]
\bar X \arrow[d,no head] \& \bar X|U \arrow[d,no head] \& \xi \arrow[l,no head] \\
S \& U \arrow[l,hook'] \&
\end{tikzcd}}
Lorsque l'on se donne une section $g$ de $\bar X$ sur $S$ contenue dans l'ouvert de lissité, on trouve de même un hom
\[
I\mathrel{\text{---}}\operatorname{Aut\,ext}(\pi_1^{(p')}(X_{\bar\eta},\xi))
\]
Ceci posé, l'image de l'un et l'autre hom. est un \uline{groupe \unsure{abélien} commutatif premier à $p^{\unsure{=|'}}$}.

Si $\bar X$ est de dim relative 1 sur $S$, alors \deleted{\unsure{la fin}} rang \inserted{topologique} du groupe image est $\leq\unsure{\nu|v}$, où $\nu$ est le nb de pts singuliers de $X_{\bar s}$.
